\subsection{Временные методы анализа EEG}

Помимо пространственной и спектральной структуры, электроэнцефалографический сигнал характеризуется выраженной временной динамикой. В задачах декодирования когнитивных состояний важно учитывать не только мгновенные характеристики сигнала, но и их изменение во времени.

\subsubsection*{Анализ отдельных проб}

Традиционный подход к анализу EEG в когнитивных исследованиях основывался на усреднении сигналов по множеству повторений стимула с целью выделения событийно-связанных потенциалов (event-related potentials, ERP). Однако в задачах машинного обучения, особенно в контексте интерфейсов мозг–компьютер (brain–computer interface, BCI), всё большее распространение получает анализ отдельных проб (single-trial analysis). Такой подход позволяет сохранять вариативность сигнала и использовать более богатую информацию о динамике состояния мозга~\footcite{10.1016/j.jneumeth.2003.10.009}.

\subsubsection*{Скользящие окна}

Одним из распространённых методов учёта временной структуры является использование скользящих окон (sliding windows). Сигнал разбивается на последовательные перекрывающиеся сегменты фиксированной длины, для каждого из которых извлекаются признаки и выполняется классификация. Это позволяет:

\begin{itemize}
\item повысить устойчивость оценки за счёт агрегации предсказаний;
\item анализировать изменение информативности во времени;
\item выявлять временные интервалы, наиболее значимые для декодирования.
\end{itemize}

Применение временной агрегации предсказаний, например на основе анализа наиболее длинной последовательности повторяющихся меток, показало способность повышать устойчивость классификации моторного воображения~\footcite{10.1109/jsen.2023.3270281}. Подобные подходы особенно актуальны в условиях низкого отношения сигнал/шум и межсубъектной вариативности.

\subsubsection*{Временно-частотные представления}

Для более детального анализа динамики спектральных характеристик используются временно-частотные представления (time-frequency representations), получаемые с помощью вейвлет-преобразования или короткого преобразования Фурье. Такие представления позволяют отслеживать изменение мощности в различных диапазонах во времени и выявлять моменты событийно-связанной десинхронизации и синхронизации~\footcite{10.1016/s1388-2457-99-00141-8}. Пример такого представления приведён на рис. \ref{fig:representing-spectral-power}.

\begin{figure}[!t]
  \centering
  \includegraphics[width=150mm]{representing-spectral-power.jpg}
  \caption{Пример представления спектральной мощности EEG во времени: исходный временной сигнал преобразуется во временно-частотную карту, где горизонтальная ось соответствует времени, вертикальная ось частоте, а цвет отражает мощность в соответствующем диапазоне.}
  \label{fig:representing-spectral-power}
\end{figure}

\subsection*{Выводы для задач реконструкции визуальных стимулов}

Рассмотренные методы предобработки и извлечения признаков формируют основу для построения систем декодирования EEG. В контексте реконструкции визуальных стимулов из сигналов, зарегистрированных в фазе воспоминания, можно выделить несколько принципиальных положений.

Во-первых, корректная стандартизированная предобработка необходима для минимизации влияния артефактов и обеспечения воспроизводимости результатов. Робастное реферирование и автоматическое выявление плохих каналов особенно важны при анализе многосубъектных данных.

Во-вторых, пространственные методы, включая CSP и римановы подходы, позволяют эффективно учитывать многоканальную структуру сигнала и снижать размерность пространства признаков. Использование ковариационных матриц и их анализа в рамках римановой геометрии представляется перспективным при ограниченном объёме обучающих данных и выраженной межсубъектной вариативности.

В-третьих, спектральные признаки, отражающие изменения мощности в альфа- и бета-диапазонах, имеют нейрофизиологическое обоснование в контексте процессов внимания и памяти, что делает их релевантными для задач визуального воображения.

Наконец, учёт временной динамики посредством скользящих окон и временной агрегации может быть полезен при низком отношении сигнал/шум, однако его эффект необходимо проверять отдельно для конкретной задачи и протокола оценки.

Таким образом, комбинирование пространственных, спектральных и временных признаков создаёт теоретическую и методологическую основу для построения моделей машинного обучения, направленных на реконструкцию структуры визуального образа по данным EEG, при условии строгого контроля разбиений и сравнения с простыми базовыми моделями.
